\documentclass[a4paper, serif, 10pt]{moderncv}

%% moderncv
% style and color
\moderncvstyle{banking}
\moderncvcolor{black}

%% page layout/margins
\usepackage[top=2.5cm, bottom=2.5cm, left=1.5cm, right=1.5cm]{geometry}

\nopagenumbers{}

%% Babel config for Spanish language
% \usepackage[spanish]{babel} has some conflicting issues with moderncv
% so we have to use enable the following options to make it work
%
% ref:
% - https://tex.stackexchange.com/a/140161/36007
\usepackage[spanish,es-lcroman]{babel}

%% fontspec
\usepackage{fontspec}

\IfFontExistsTF{Linux Libertine}{
  \setmainfont[Ligatures={TeX, Common}, Numbers=OldStyle]{Linux Libertine}
}{}
\IfFontExistsTF{Linux Libertine O}{
  \setmainfont[Ligatures={TeX, Common}, Numbers=OldStyle]{Linux Libertine O}
}{}

%% CTeX
% CJK support, used to show CJK characters in the resume
%
% - fontset=none: disable builtin fontset but instead set the CJK font manually
% - heading=false: disable ctex heading
% - punct=kaiming: use kaiming punctuations styles for CJK
% - scheme=plain: use plain scheme, do not override \normalsize font size
% - space=auto: space settings for CJK characters
%
% ref:
% - http://ctan.mirrorcatalogs.com/language/chinese/ctex/ctex.pdf
\usepackage[UTF8, heading=false, punct=kaiming, scheme=plain, space=auto]{ctex}

\IfFontExistsTF{Noto Serif CJK SC}{
  \setCJKmainfont{Noto Serif CJK SC}
}{}
\IfFontExistsTF{Noto Sans CJK SC}{
  \setCJKsansfont{Noto Sans CJK SC}
}{}

%% line spacing
\usepackage{setspace}
\setstretch{1.125}

%% URL styling - use normal text instead of monospace
\urlstyle{same}

% Auto-underline all links
% Defer until \begin{document} so hyperref is already loaded
\AtBeginDocument{%
  \let\oldhref\href
  \renewcommand{\href}[2]{\oldhref{#1}{\underline{#2}}}%
}

%% Patch \httplink and \httpslink to handle full URLs
% The original moderncv macros always prepend http:// or https:// to URLs.
% This causes issues when the URL already has a protocol (e.g., https://example.com)
% resulting in malformed URLs like https://https://example.com.
% This patch checks if the URL already contains "://" and uses it directly if so.
% Note: We use \str_set:Nx (x-expansion) to expand macros like \@homepage first.
\ExplSyntaxOn
\str_new:N \g_yamlresume_url_str

\RenewDocumentCommand{\httpslink}{O{}m}{
  \str_set:Nx \g_yamlresume_url_str {#2}
  \str_if_in:NnTF \g_yamlresume_url_str {://}
    { \url{#2} }
    { \url{https://#2} }
}

\RenewDocumentCommand{\httplink}{O{}m}{
  \str_set:Nx \g_yamlresume_url_str {#2}
  \str_if_in:NnTF \g_yamlresume_url_str {://}
    { \url{#2} }
    { \url{http://#2} }
}
\ExplSyntaxOff

\name{Lucía Fernández Gómez}{}
\title{Customer Success Manager | Especialista en Retención y Crecimiento de Clientes}
\phone[mobile]{+34 612 345 678}
\email{lucia.fernandez@example.com}
\homepage{https://www.luciafernandez.es}

\address{Calle de Serrano 45, Madrid, Comunidad de Madrid, España, 28001}{}{}

\extrainfo{{\small \faLinkedin}\ \href{https://www.linkedin.com/in/lucia-fernandez-csm}{@lucia-fernandez-csm} {} {} {} • {} {} {} 
{\small \faTwitter}\ \href{https://twitter.com/lucia\_csm}{@lucia\_csm}}

\begin{document}

\maketitle

\section{Información básica}

\cvline{}{Customer Success Manager con más de 8 años de experiencia en empresas SaaS y tecnología, enfocada en la retención, adopción y crecimiento de clientes estratégicos.
%
\begin{itemize}
\item Gestión de carteras de clientes B2B con un NPS superior a 70 y una tasa de retención del 95\%.
\item Implementación de programas de onboarding que reducen el time-to-value en un 30\%.
\item Colaboración estrecha con Ventas, Producto y Soporte para identificar oportunidades de upselling y cross-selling.
\item Comunicación bilingüe (español nativo, inglés profesional) y orientación a resultados.
\end{itemize}}
\section{Educación}

\cventry{sept 2012–jul 2016}
        {Licenciatura, Administración y Dirección de Empresas, Puntuación: 8.5}
        {Universidad Complutense de Madrid}
        {\url{https://www.ucm.es}}
        {}
        {\begin{itemize}
\item Formación sólida en gestión empresarial, análisis financiero y estrategia comercial.
\item Participación en programas de intercambio y proyectos de consultoría para pymes.
\item Desarrollo de habilidades de comunicación, liderazgo y trabajo en equipo.
\end{itemize}
\textbf{Cursos}: Fundamentos de Marketing, Comportamiento Organizacional, Gestión de Recursos Humanos, Contabilidad Financiera, Estadística Aplicada a la Empresa, Sistemas de Información para la Gestión, Negociación y Habilidades Directivas, Estrategia Empresarial}
\section{Experiencia laboral}

\cventry{ene 2021 hasta la fecha}
        {Senior Customer Success Manager}
        {Nexo Cloud Solutions}
        {\url{https://www.nexocloud.es}}
        {}
        {\begin{itemize}
\item Responsable de una cartera de 45 clientes enterprise con un valor anual contratado (ACV) de 4,5 M€.
\item Diseñé y ejecuté el plan de éxito del cliente, logrando un aumento del 25\% en la adopción de producto y una reducción del churn del 18\%.
\item Lideré revisiones trimestrales (QBR) con directivos, alineando métricas de negocio con el valor entregado.
\item Identifiqué oportunidades de expansión por 1,2 M€ mediante análisis de uso y colaboración con el equipo de Ventas.
\item Mentoría a dos Customer Success Managers junior, mejorando sus índices de satisfacción y retención.
\end{itemize}
\textbf{Palabras clave}: Retención de clientes, Onboarding, QBR, Upselling, SaaS, CRM}

\cventry{mar 2018–dic 2020}
        {Customer Success Manager}
        {DataVista Analytics}
        {\url{https://www.datavista.es}}
        {}
        {\begin{itemize}
\item Gestioné una cartera de 80 clientes pyme y mid-market en el sector de analítica de datos.
\item Implementé un programa de onboarding automatizado que redujo el tiempo de activación de 30 a 21 días.
\item Aumenté el NPS de 45 a 68 en 18 meses mediante seguimiento proactivo y resolución de incidencias.
\item Colaboré con Producto en la priorización de funcionalidades solicitadas por clientes clave.
\item Realicé formaciones periódicas y webinars para fomentar la adopción de nuevas funcionalidades.
\end{itemize}
\textbf{Palabras clave}: Analítica de datos, NPS, Onboarding, Formación, Pyme}

\cventry{sept 2016–feb 2018}
        {Customer Support Specialist}
        {Soluciones Ágiles}
        {\url{https://www.solucionesagiles.com}}
        {}
        {\begin{itemize}
\item Atendí solicitudes de soporte técnico y funcional de más de 200 clientes corporativos.
\item Resolví incidencias con un índice de satisfacción del 92\% y un tiempo medio de respuesta inferior a 2 horas.
\item Documenté procedimientos y creé una base de conocimiento que redujo las consultas repetitivas en un 35\%.
\item Coordiné escalados con los equipos de Ingeniería y Producto para la resolución de bugs críticos.
\end{itemize}
\textbf{Palabras clave}: Soporte técnico, Base de conocimiento, SLA, Escalado}
\section{Idiomas}

\cvline{Español}{Competencia nativa o bilingüe \hfill \textbf{Palabras clave}: Español nativo}
\cvline{Inglés}{Competencia profesional plena \hfill \textbf{Palabras clave}: C1, TOEIC 950}
\cvline{Francés}{Competencia elemental \hfill \textbf{Palabras clave}: A2}
\section{Competencias}

\cvline{Gestión de Customer Success}{Experto \hfill \textbf{Palabras clave}: Retención, Onboarding, QBR, Planes de éxito, Renovaciones}
\cvline{Herramientas CRM}{Avanzado \hfill \textbf{Palabras clave}: Salesforce, HubSpot, Gainsight, Zendesk}
\cvline{Análisis de Datos}{Avanzado \hfill \textbf{Palabras clave}: Excel, SQL, Tableau, KPIs}
\cvline{Comunicación y Liderazgo}{Experto \hfill \textbf{Palabras clave}: Presentaciones, Negociación, Mentoría, Trabajo en equipo}
\section{Premios}

\cventry{dic 2023}
        {Nexo Cloud Solutions}
        {Premio a la Excelencia en Customer Success}
        {}
        {}
        {Reconocimiento anual por lograr la mayor tasa de retención y satisfacción de clientes en la compañía.}

\cventry{jun 2020}
        {DataVista Analytics}
        {Innovación en Procesos de Onboarding}
        {}
        {}
        {Galardón por el diseño e implementación de un programa de onboarding que mejoró el time-to-value en un 30\%.}
\section{Certificados}

\cventry{mar 2021}
        {SuccessCOACHING}
        {Certified Customer Success Manager (CCSM)}
        {\url{https://www.successcoaching.com/certifications}}
        {}
        {}

\cventry{sept 2019}
        {Salesforce}
        {Salesforce Certified Administrator}
        {\url{https://trailhead.salesforce.com/certifications}}
        {}
        {}
\section{Publicaciones}

\cventry{may 2022}
        {Estrategias de retención en SaaS B2B}
        {Revista Customer Success España}
        {\url{https://www.revistacustomersuccess.es/articulos/retencion-saas-b2b}}
        {}
        {\begin{itemize}
\item Artículo sobre las mejores prácticas para reducir el churn en empresas de software como servicio.
\item Incluye casos de éxito y métricas clave para medir el impacto de Customer Success.
\end{itemize}}
\section{Proyectos}

\cventry{ene 2022–jun 2022}
        {Rediseño del proceso de incorporación de clientes para acelerar la adopción y el valor.}
        {Programa de Onboarding 2.0}
        {\url{https://www.nexocloud.es/onboarding}}
        {}
        {\begin{itemize}
\item Reestructuré las fases de onboarding con hitos medibles y materiales personalizados por sector.
\item Integré el programa con el CRM para automatizar seguimientos y alertas tempranas.
\item Resultado: reducción del tiempo de activación en un 30\% y aumento del 20\% en la adopción de funciones clave.
\end{itemize}
\textbf{Palabras clave}: Onboarding, Automatización, Adopción}

\cventry{mar 2021–sept 2021}
        {Panel interno para monitorizar el riesgo de churn y las oportunidades de expansión.}
        {Dashboard de Salud del Cliente}
        {\url{https://www.nexocloud.es/dashboard-salud}}
        {}
        {\begin{itemize}
\item Definí métricas de salud del cliente (uso, tickets, NPS, facturación) y umbrales de alerta.
\item Desarrollé informes automatizados en Tableau para el equipo de Customer Success.
\item Permitiió priorizar acciones proactivas y evitar renovaciones en riesgo.
\end{itemize}
\textbf{Palabras clave}: Churn, KPIs, Tableau}
\section{Intereses}

\cvline{Tecnología}{SaaS, Inteligencia artificial, No-code}
\cvline{Sostenibilidad}{Economía circular, Voluntariado ambiental}
\cvline{Bienestar}{Yoga, Meditación, Running}
\section{Voluntariado}

\cventry{sept 2021 hasta la fecha}
        {Mentor de Customer Success}
        {Asociación Española de Customer Success}
        {\url{https://www.aecs.es}}
        {}
        {\begin{itemize}
\item Ejercí como mentora de profesionales junior en la definición de planes de carrera y mejora de habilidades de gestión de clientes.
\item Participé en jornadas y talleres sobre retención, onboarding y métricas de éxito del cliente.
\item Colaboré en la creación de una guía de buenas prácticas para equipos de Customer Success en España.
\end{itemize}}
    
\end{document}